\documentclass[11pt]{article}
\usepackage[T1]{fontenc}
\usepackage[margin=0.88in]{geometry}
\usepackage{amsmath,amssymb,amsthm}
\usepackage{booktabs,tabularx}
\usepackage[hidelinks]{hyperref}
\usepackage{microtype}

\newtheorem{theorem}{Theorem}
\newcommand{\N}{\mathbb N}
\newcommand{\FHC}{\mathrm{FHC}}

\title{Finite self-direct-sums for a regular C-type\\
frequently hypercyclic class}
\author{New partial result for a question in arXiv:2007.12647}
\date{11 August 2026}

\begin{document}
\maketitle

\paragraph{Status.}
\textbf{New partial result; the unrestricted question remains open.}
For the principal regular \(C_{+,1}\)-type class, including the known chaotic
frequently hypercyclic example which is not hereditarily frequently
hypercyclic, frequent hypercyclicity survives every finite self-direct-sum.
Thus that central candidate family cannot yield a counterexample to the
question whether \(T\oplus T\) is frequently hypercyclic whenever \(T\) is.

\section*{Source question and later boundary}
Gilmore's survey \cite{survey} asks whether frequent hypercyclicity of a
bounded operator \(T\) implies frequent hypercyclicity of \(T\oplus T\).
Here \(T\) is frequently hypercyclic if some vector visits every nonempty open
set at a set of times of positive lower density.

The question is still stated as open in \cite{hfhc}.  That paper introduces
hereditary frequent hypercyclicity, a stronger property which does imply the
desired conclusion, and then constructs a chaotic frequently hypercyclic
\(C_{+,1}\)-type operator which is not hereditary.  The construction therefore
tests whether failure of heredity might also break the self-direct-sum.

\section*{The C-type parameters}
We use the notation of \cite{gmm}.  At level \(k\), the relevant
\(C_{+,1}\)-type blocks have common length \(\Delta_k\), feedback magnitude
\(2^{-\tau_k}\), and internal weights
\[
 w_i^{(k)}=
 \begin{cases}
 2,&1\le i\le\delta_k,\\
 1,&\delta_k<i<\Delta_k.
 \end{cases}
\]
The block map sends each sufficiently high block to an earlier block, and
every finitely supported vector is periodic.  The quantity
\[
 \eta:=\limsup_{k\to\infty}
 \frac{\delta_k-\tau_k}{\Delta_k}
\]
measures the positive amplification margin available for periodic orbit
shadowing.

\begin{theorem}[finite-self-sum theorem]\label{thm:main}
Let \(1\le p<\infty\), and let \(T\) be a \(C_{+,1}\)-type operator on
\(\ell_p(\N)\) with \(\eta>0\).  Then, for every integer \(r\ge1\),
\[
 T^{\oplus r}:=\underbrace{T\oplus\cdots\oplus T}_{r\text{ copies}}
\]
is frequently hypercyclic.
\end{theorem}

The same proof works under the more general pair of weight-product hypotheses
in the periodic-shadowing theorem of \cite{gmm}; the displayed limsup is a
convenient readable sufficient condition.

\section*{Idea of the proof}
The one-copy proof in \cite[Thms.~5.31 and 6.9]{gmm} does more than produce one frequently
hypercyclic vector.  Given a finitely supported periodic target \(x\), it
chooses one high block level \(k\), one transfer time \(m\), and one common
period \(d=2\Delta_k\).  A correction \(z=L_{k,m}x\), linear in \(x\), is
small along a fixed positive fraction of its period and, after time \(m\),
shadows the same fraction of the orbit of \(x\).

For finitely many targets \(x_1,\ldots,x_r\), choose a level above all their
supports and make the same weight-product lower bound large enough for the
largest of their norms.  The same \(k,m,d\) then works in every coordinate.
The tuple \((L_{k,m}x_1,\ldots,L_{k,m}x_r)\) supplies exactly the periodic
shadowing condition for \(T^{\oplus r}\).

\section*{Proof}
Let \(X_0=c_{00}\subset\ell_p(\N)\).  It is dense, \(T\)-invariant, and
consists of periodic vectors.  The frequent-hypercyclicity criterion used in
\cite{gmm} has the following operative form.  For some \(\rho>0\), for every
\(x\in X_0\), tolerance \(\varepsilon>0\), and previously fixed finite
periodic data, one finds a periodic \(z\in X_0\) and integers \(m,d\ge1\)
such that \(d\) is a common admissible period and
\begin{equation}\label{eq:shadow}
 \|T^jz\|<\varepsilon,
 \qquad
 \|T^{m+j}z-T^jx\|<\varepsilon
 \quad(0\le j\le \rho d).
\end{equation}
This condition implies frequent hypercyclicity by concatenating the positive
proportions of periodic shadowing times.

For \(C_{+,1}\)-type operators, the proof of the periodic-shadowing theorem
constructs \(z\) in descendant blocks at a level \(k\).  For fixed \(k,m\),
its formula (displayed in the proof of Theorem~6.7, to which the proof of
Theorem~6.9 explicitly refers) is linear in the coefficients of \(x\).  If
the support of \(x\)
lies below level \(k_0\), both errors in \eqref{eq:shadow} are bounded by
\[
 C^{-1}A^{\Delta_{k_0}}\|x\|,
\]
where \(A\) is independent of \(x\), once the two level-\(k\) weight products
are at least \(C\).  The correction has common period \(d=2\Delta_k\), and
the source's divisibility choice accommodates the earlier periodic data.

We verify that arbitrarily large common bounds \(C\) are available.  Choose
\(0<a<\eta/2\).  For arbitrarily large \(k\),
\(\delta_k-\tau_k>2a\Delta_k\).  Taking \(k\) larger still makes
\(a\Delta_k>\log_2 C\).  With the source proof's choice \(m=\Delta_k-1\),
the first relevant product is
\[
 2^{-\tau_k}\prod_{i=1}^{\Delta_k-1}w_i^{(k)}
 =2^{\delta_k-\tau_k}>C,
\]
and, uniformly for \(0\le j\le a\Delta_k\), the second is at least
\[
 2^{\delta_k-\tau_k-j}>2^{a\Delta_k}>C.
\]
These are precisely the amplification bounds used to obtain
\eqref{eq:shadow}: since \(d=2\Delta_k\), take \(\rho=a/2\), so that
\(\rho d=a\Delta_k\).

Now fix \(r\), use the max norm on the finite product, and take
\(\mathbf x=(x_1,\ldots,x_r)\in X_0^r\).  Choose one \(k_0\) above all
supports and choose \(C\) so large that
\(C^{-1}A^{\Delta_{k_0}}\max_i\|x_i\|<\varepsilon\).
Use the same good \(k,m\) in all coordinates and put
\[
 \mathbf z=(L_{k,m}x_1,\ldots,L_{k,m}x_r).
\]
Every coordinate has the same admissible period \(d\), and the coordinatewise
estimates give \eqref{eq:shadow} for \(T^{\oplus r}\) and
\(\mathbf x\).  Since \(X_0^r\) is dense, invariant, and periodic, the same
criterion proves that \(T^{\oplus r}\) is frequently hypercyclic.

\section*{Application to the non-hereditary example}
In the example of \cite{hfhc},
\[
 \delta_k=\frac{\Delta_k}{4},
 \qquad
 \tau_k=\frac{\Delta_k}{8},
 \qquad
 \frac{\delta_k-\tau_k}{\Delta_k}=\frac18.
\]
The block lengths are then chosen sufficiently rapidly to prove that \(T\)
is not frequently hypercyclic along one positive-lower-density set.  Hence
\(T\) is chaotic and frequently hypercyclic but not hereditarily frequently
hypercyclic.  Theorem~\ref{thm:main} nevertheless gives
\(T^{\oplus r}\in\FHC\) for every finite \(r\).

There is also a broader regular-class corollary.  Under the monotonicity and
summability hypotheses in the \(C_{+,1}\) equivalence theorem of \cite{gmm},
and assuming
\(\limsup\tau_k/\delta_k<1\), frequent hypercyclicity is equivalent to
\(\limsup\delta_k/\Delta_k>0\).  It therefore forces \(\eta>0\), so every
frequently hypercyclic operator in that regular class has frequently
hypercyclic finite self-sums.

\section*{Scope and failed full upgrades}
\begin{center}
\small
\begin{tabularx}{\textwidth}{@{}p{0.32\textwidth}X@{}}
\toprule
Route & Outcome \\
\midrule
Regular \(C_{+,1}\) candidates & Ruled out by the theorem above. \\
Single bounded-below weighted shifts & Known finite-self-sum stability; no counterexample. \\
Two-shift bad pair encoded into one shift & The incompatible support mechanism vanishes for identical copies. \\
Common extension of the bad pair & No frequently hypercyclic extension with the two required quotient maps was found. \\
Invariant-measure products & Frequent hypercyclicity alone supplies neither a full-support weakly mixing measure nor an adequate product theorem. \\
Arbitrary C-type operators & No positive shadowing margin follows from frequent hypercyclicity outside the regular subclass. \\
\bottomrule
\end{tabularx}
\end{center}

Thus this is a substantive elimination theorem, not a solution of the
unrestricted question.  A full counterexample must leave both the standard
weighted-shift mechanism and the synchronized periodic-shadowing \(C\)-type
regime, or find a new way to defeat simultaneous positive lower densities.

\paragraph{Novelty and verification.}
Exact arXiv-id, phrase, and core-keyword searches through 11 August 2026 found
the survey question, the 2024 still-open statement, and the two-different-shift
counterexample, but no prior statement of the finite-self-sum consequence
above.  The proof is an explicit coordinatewise audit of the published
periodic-shadowing estimates.  Human review should focus on the common-period
divisibility step and the uniformity of the constants in the transfer formula.

\begin{thebibliography}{9}
\raggedright
\bibitem{survey}
E. Gilmore,
\emph{Linear Dynamical Systems}, arXiv:2007.12647.

\bibitem{gmm}
S. Grivaux, E. Matheron, and Q. Menet,
\emph{Linear dynamical systems on Hilbert spaces: typical properties and
explicit examples}, arXiv:1703.01854.  See Theorems~5.31, 6.9, 6.17,
and 6.18, together with the transfer formula in the proof of Theorem~6.7.

\bibitem{hfhc}
F. Bayart, S. Grivaux, E. Matheron, and Q. Menet,
\emph{Hereditarily frequently hypercyclic operators and disjoint frequent
hypercyclicity}, arXiv:2409.07103.  See the still-open self-direct-sum question
on p.~2 and the chaotic non-hereditarily frequently hypercyclic
\(C_{+,1}\) example on p.~25.
\end{thebibliography}

\end{document}
